\section{Options Protocol}


\subsection{European vs American Options}

We support both European (exercise at expiry) and American (exercise anytime) options:

\begin{definition}[Option Contract]
An option $O$ is characterized by:
\begin{equation}
O = \{S, K, T, \sigma, r, \phi, \text{style}\}
\end{equation}
where $S$ is spot price, $K$ is strike, $T$ is time to expiry, $\sigma$ is implied volatility, $r$ is risk-free rate, $\phi \in \{\text{call}, \text{put}\}$, and style $\in \{\text{European}, \text{American}\}$.
\end{definition}

\subsection{Black-Scholes Pricing}

For European options, we use Black-Scholes formula:

\begin{align}
C &= S \cdot N(d_1) - K \cdot e^{-rT} \cdot N(d_2) \\
P &= K \cdot e^{-rT} \cdot N(-d_2) - S \cdot N(-d_1)
\end{align}

where:
\begin{align}
d_1 &= \frac{\ln(S/K) + (r + \sigma^2/2)T}{\sigma\sqrt{T}} \\
d_2 &= d_1 - \sigma\sqrt{T}
\end{align}

\begin{lstlisting}[caption={Black-Scholes Implementation}]
library BlackScholes {
    using FixedPoint for uint256;

    function calculateCallPrice(
        uint256 S,     // Spot price
        uint256 K,     // Strike price
        uint256 T,     // Time to expiry (in years)
        uint256 sigma, // Implied volatility
        uint256 r      // Risk-free rate
    ) public pure returns (uint256) {
        uint256 d1 = calculateD1(S, K, T, sigma, r);
        uint256 d2 = d1.sub(sigma.mul(sqrt(T)));

        uint256 Nd1 = normalCDF(d1);
        uint256 Nd2 = normalCDF(d2);

        uint256 discountFactor = exp(r.mul(T).neg());

        return S.mul(Nd1).sub(K.mul(discountFactor).mul(Nd2));
    }

    function calculateD1(
        uint256 S,
        uint256 K,
        uint256 T,
        uint256 sigma,
        uint256 r
    ) private pure returns (uint256) {
        uint256 numerator = ln(S.div(K)).add(
            r.add(sigma.pow(2).div(2)).mul(T)
        );
        uint256 denominator = sigma.mul(sqrt(T));
        return numerator.div(denominator);
    }
}
\end{lstlisting}

\subsection{Implied Volatility Calculation}

We calculate IV using Newton-Raphson iteration:

\begin{equation}
\sigma_{n+1} = \sigma_n - \frac{V_{\text{BS}}(\sigma_n) - V_{\text{market}}}{\text{vega}(\sigma_n)}
\end{equation}

Convergence criterion: $|\sigma_{n+1} - \sigma_n| < 10^{-6}$.

\subsection{Greeks Calculation}

Option sensitivities guide risk management:

\subsubsection{Delta ($\Delta$)}
Price sensitivity to underlying:
\begin{align}
\Delta_{\text{call}} &= N(d_1) \\
\Delta_{\text{put}} &= N(d_1) - 1
\end{align}

\subsubsection{Gamma ($\Gamma$)}
Delta sensitivity to underlying:
\begin{equation}
\Gamma = \frac{\phi(d_1)}{S \cdot \sigma \cdot \sqrt{T}}
\end{equation}

\subsubsection{Theta ($\Theta$)}
Time decay:
\begin{align}
\Theta_{\text{call}} &= -\frac{S \cdot \phi(d_1) \cdot \sigma}{2\sqrt{T}} - r \cdot K \cdot e^{-rT} \cdot N(d_2) \\
\Theta_{\text{put}} &= -\frac{S \cdot \phi(d_1) \cdot \sigma}{2\sqrt{T}} + r \cdot K \cdot e^{-rT} \cdot N(-d_2)
\end{align}

\subsubsection{Vega ($\nu$)}
Volatility sensitivity:
\begin{equation}
\nu = S \cdot \phi(d_1) \cdot \sqrt{T}
\end{equation}

\subsubsection{Rho ($\rho$)}
Interest rate sensitivity:
\begin{align}
\rho_{\text{call}} &= K \cdot T \cdot e^{-rT} \cdot N(d_2) \\
\rho_{\text{put}} &= -K \cdot T \cdot e^{-rT} \cdot N(-d_2)
\end{align}

